\newpage
\section{Misión 3: Pastoreando el caos}\label{sec:mission3}

\begin{minipage}{0.45\linewidth}
Debido a diferencias políticas, y a pesar de los esfuerzos de la asamblea de la colonia, cada base del planeta adoptó una tecnología de LAN (red de área local) distinta (!).
Aunque las comunicaciones son satisfactorias dentro de cada base, por el momento es imposible conectar directamente dos ordenadores cualesquiera de dos bases diferentes.
%
Una vez más, la colonia te pide ayuda.\\
\end{minipage}
\hspace{0.04\linewidth}
\begin{minipage}{0.5\linewidth}
\vspace{-1cm}\includegraphics[width=\linewidth]{colony_network_map.pdf}
\end{minipage}

\noindent
Gracias a los sacrificios de todos, la colonia está ahora en condiciones de construir e instalar un nuevo router (encaminador) y dos buses cableados.
Estos se representan con una línea más gruesa en el diagrama de abajo.
%
% \noindent
Tu objetivo a largo plazo es hacer posible la comunicación entre redes entre dos nodos cualesquiera de dos bases cualesquiera.
Por ahora, diseñarás la infraestructura física y lógica.\\

\subsection*{Detalles del informe}

Todas las bases emplean el protocolo IPv4, y hay varias bases conectadas tanto al bus Norte como al bus Sur.
Cada base contiene una única LAN. Coordínate con los demás para que \textit{no haya dos LANs del mismo tamaño}.

\begin{enumerate}[itemsep=0.1cm]
\item Dentro de tu grupo:
\begin{itemize}[itemsep=0.1cm]
    \item Determina cuántas direcciones necesitas para que cada dispositivo pueda identificarse de forma única dentro de tu LAN.
    Elige el tamaño mínimo de LAN que puede satisfacer este requisito, y cuantifica el ``espacio'' libre disponible en
    tu LAN.
    \item Indica cuántos de los 32~bits de una dirección IPv4 necesitas para identificar un único dispositivo conectado a la LAN de tu base.
    ¿Qué concepto usado en IP está más relacionado con este número? ¿Cómo puede ser útil esta estructura?\\[0.2cm]
\end{itemize}

\begin{center}
\begin{bytefield}{18}
\bitbox{8}{Network ID} & \bitbox{15}{Host ID}\\
\end{bytefield}
\end{center}

\item En coordinación con los otros grupos del aula:
\begin{itemize}[itemsep=0.1cm]
\item Crea una LAN para el bus Norte, y describe todas las direcciones físicas y lógicas conectadas a ella.
\item Asegúrate de que los bloques de direcciones IPv4 asignados a las LANs son de tamaños diferentes y no se solapan.
\item Usando la pizarra o un documento compartido, enumera todas las LANs, sus tamaños y sus rangos de direcciones (actualizados).
\end{itemize}

\end{enumerate}

\subsection*{Detalles del feedback}

Utiliza un proceso de feedback similar al de las misiones anteriores, asegurándote de tratar:

\begin{enumerate}[itemsep=0.1cm]
\item ¿Están las LANs correctamente dimensionadas?
\item ¿Hay alguna colisión entre los bloques?
\item ¿Extrajo el informe original toda la información sobre la estructura de la dirección IP?
\end{enumerate}

